\section{Hạn chế, nguy cơ đối với tính hợp lệ và sử dụng có trách nhiệm}
\label{sec:limitations}

Mặc dù thực nghiệm được thiết kế theo dạng paired comparison và toàn bộ
$250$ mẫu đều được xử lý thành công, kết quả cần được diễn giải trong phạm
vi của corpus, benchmark và implementation hiện tại. Phần này phân tích các
nguy cơ đối với tính hợp lệ cấu trúc, nội tại, ngoại tại và kết luận, đồng
thời xác định giới hạn sử dụng của hệ thống.

\subsection{Tính hợp lệ cấu trúc}

\paragraph{Quan hệ điều kiện--hệ quả không đồng nhất với quan hệ nhân quả
thực nghiệm.}

Đồ thị được xây dựng bằng cách thao tác hóa các quy tắc pháp luật thành quan
hệ điều kiện--hệ quả. Các cạnh \texttt{CAUSES} biểu diễn cơ chế quy phạm đã
được mã hóa, không phải quan hệ nhân quả được khám phá từ dữ liệu quan sát
hoặc ước lượng từ thí nghiệm. Vì vậy, LegalSCM chỉ là mô hình nhân quả cấu
trúc pháp lý trong phạm vi biểu diễn quy tắc đã cho. Kết quả của phép
$\doop(M=m)$ là kết quả của mô phỏng trên mô hình này, không phải ước lượng
tác động nhân quả trong thế giới thực \cite{pearl2009causality}.

\paragraph{Counterfactual subset chưa đo đầy đủ hard intervention.}

Hai mươi mẫu được gắn cờ \texttt{requires\_counterfactual=true} chủ yếu là
các direct-edge negative: graph chứa đường $A\rightarrow M\rightarrow Y$
nhưng không chứa cạnh trực tiếp $A\rightarrow Y$. Chúng đo khả năng phân
biệt quan hệ trực tiếp và gián tiếp tốt hơn là khả năng dự đoán một
counterfactual world sau can thiệp. Do đó, Counterfactual Accuracy
$100\%$ không xác nhận tính đúng đắn tổng quát của structural intervention.

\paragraph{Gold path tuyến tính không bao phủ cấu trúc phi tuyến.}

Trong $250$ mẫu, chỉ $236$ mẫu có gold path tuyến tính. Bốn mẫu branch và
mười mẫu convergence không có một chuỗi duy nhất biểu diễn đầy đủ cấu trúc
cần trả lời, nhưng evaluator hiện vẫn chấm bằng exact linear path. Vì vậy,
Top-1 Exact Path và Oracle Exact Path bằng $0\%$ ở hai nhóm này có thể phản
ánh giới hạn của metric thay vì lỗi hoàn toàn của hệ thống.

\paragraph{Metric câu trả lời và citation chỉ là proxy.}

Token F1 và ROUGE-L đo mức chồng lấp bề mặt, không bảo đảm một câu trả lời
đúng về mặt pháp lý. Citation metric so sánh tập định danh điều luật nhưng
chưa kiểm tra semantic entailment giữa từng assertion và điều luật được
trích dẫn. Normalized Exact Match cũng không phù hợp với câu trả lời có thêm
giải thích và marker evidence. Vì vậy, các metric tự động không thể thay thế
đánh giá của chuyên gia.

\subsection{Tính hợp lệ nội tại}

\paragraph{Tính vòng tròn giữa graph và benchmark.}

Câu hỏi, gold rule, gold event và nhiều negative được xây dựng từ cùng graph
mà hệ thống truy hồi. Thiết kế này hữu ích để kiểm tra khả năng khôi phục một
encoding đã biết, nhưng có thể làm tăng độ tương thích giữa benchmark và
biểu diễn của pipeline. Kết quả không chứng minh hệ thống có thể xử lý một
câu hỏi pháp luật được tạo độc lập hoặc một cách diễn đạt ngoài phân bố.

\paragraph{Paired ablation được kiểm soát nhưng chưa đủ phân biệt.}

Hai cấu hình dùng cùng benchmark, retrieval resource, threshold, answer
generator và evaluator; biến thay đổi là
\texttt{counterfactual\_mode}. Điều này giảm confound trong so sánh. Tuy
nhiên, benchmark thiếu các trường hợp có mechanism thay thế, exception,
conflict hoặc unknown root làm hard intervention và node deletion đưa ra kết
quả khác nhau. Vì thế, null quality result có thể đến từ độ nhạy thấp của
benchmark chứ không chứng minh hai phương pháp tương đương.

\paragraph{Query analysis và event matching còn mang tính heuristic.}

Việc xác định claim type, answer intent và mediator dựa một phần trên biểu
thức chính quy tiếng Việt, từ khóa và semantic mapping. Các heuristic này có
thể thất bại trước phủ định phức tạp, cách diễn đạt mới, thuật ngữ đồng âm
hoặc truy vấn chứa nhiều claim. Sai sót ở đây có thể ảnh hưởng đồng thời đến
verification decision và dạng câu trả lời.

\paragraph{Runtime chưa được tái lập trên phần cứng được công bố.}

Hai run được thực hiện trên cùng server, nhưng hệ điều hành, phiên bản
Python, CPU, RAM và GPU chưa được điền vào Bảng~\ref{tab:environment}. Thời
gian structural SCM cao hơn khoảng $3{,}23$ lần theo trung bình mỗi mẫu và
$2{,}85$ lần theo wall-clock chỉ có giá trị đối với môi trường đã chạy. Chưa
có thí nghiệm lặp với cold cache, warm cache hoặc nhiều cấu hình phần cứng.

\subsection{Hạn chế dữ liệu và mô hình}

\paragraph{Chuẩn hóa quy tắc có thể làm mất sắc thái pháp lý.}

Corpus chuyển văn bản luật thành condition, effect, subject và article.
Quá trình này có thể giản lược các biểu thức tình thái như ``có thể'', điều
kiện kết hợp, ngoại lệ, dẫn chiếu chéo và phạm vi áp dụng. Event
canonicalization cũng có thể hợp nhất hai biểu thức gần nghĩa nhưng khác
biệt về pháp lý, hoặc giữ nhiều alias trong cùng display name.

\paragraph{LegalSCM là mô hình tất định và giới hạn.}

Mô hình sử dụng ba trạng thái
$\{\mathrm{TRUE},\mathrm{FALSE},\mathrm{UNKNOWN}\}$ và fixed-point
inference trên các rule đã mã hóa. Phiên bản hiện tại chưa biểu diễn đầy đủ:

\begin{itemize}
    \item nhiều điều kiện đồng thời và các cấu trúc điều kiện lồng nhau;
    \item negation, exception và defeasible rule;
    \item thứ bậc văn bản, ưu tiên quy tắc và xung đột cách giải thích;
    \item hiệu lực theo thời gian và thay đổi của pháp luật;
    \item xác suất, mức độ tin cậy hoặc bất định nhận thức;
    \item bước abduction từ dữ kiện quan sát sang factual context.
\end{itemize}

Do đó, LegalSCM không phải mô hình hoàn chỉnh của suy luận pháp lý hoặc quan
hệ nhân quả ngoài đời thực.

\paragraph{Search bị giới hạn bởi ngân sách.}

Retriever mở rộng tối đa hai hop, trong khi path ablation tìm alternative
path tối đa ba hop và giới hạn số candidate. Phát biểu ``không tồn tại đường
thay thế'' trong artifact chỉ có nghĩa là không tìm thấy đường trong graph
và ngân sách cấu hình, không phải chứng minh không tồn tại bất kỳ cơ chế
pháp lý nào khác.

\subsection{Tính hợp lệ ngoại tại}

Corpus chỉ tập trung vào Bộ luật Hình sự Việt Nam năm 2015 và các nội dung
sửa đổi được đưa vào dữ liệu. Nghiên cứu chưa đánh giá:

\begin{itemize}
    \item văn bản pháp luật thuộc lĩnh vực dân sự, hành chính hoặc thương mại;
    \item án lệ, nghị quyết hướng dẫn, hồ sơ vụ án hoặc thực tiễn xét xử;
    \item văn bản mới hoặc cách diễn đạt không xuất hiện trong corpus;
    \item câu hỏi do người dùng thực tế hoặc chuyên gia độc lập biên soạn;
    \item ngôn ngữ khác ngoài tiếng Việt;
    \item graph có quy mô và độ sâu lớn hơn đáng kể.
\end{itemize}

Do benchmark tổng hợp được tạo từ cùng nguồn graph, không thể suy rộng trực
tiếp kết quả sang hỏi--đáp pháp luật mở hoặc sang hệ thống tư vấn pháp lý
trong thực tế.

\subsection{Tính hợp lệ kết luận}

Benchmark gồm $250$ mẫu, trong đó một số nhóm rất nhỏ, chẳng hạn branch chỉ
có bốn mẫu và bridge-convergence fallback có sáu mẫu. Các tỷ lệ theo nhóm
này có phương sai lớn và không nên được xem là ước lượng ổn định.

Nghiên cứu mới có một run hoàn chỉnh cho mỗi chế độ trên bộ artifact chính,
chưa báo cáo confidence interval hoặc kiểm định thống kê theo mẫu. Do hai
cấu hình tạo quality metric giống hệt nhau, không có cơ sở để tuyên bố
structural SCM cải thiện accuracy. Tương tự, Oracle Exact Path chỉ là phân
tích upper bound trên candidate set và không đại diện cho hiệu năng có thể
triển khai của Top-1 selector.

Difficulty label cũng chưa được hiệu chỉnh độc lập với question type. Việc
nhóm hard đạt Answer F1 cao hơn nhóm medium phản ánh thành phần mẫu, không
chứng minh hệ thống xử lý câu hỏi khó tốt hơn.

\subsection{Thẩm định chuyên gia và tính cập nhật pháp luật}

Rules, event normalization, câu hỏi, gold answer, gold path và citation chưa
được một nhóm chuyên gia pháp luật độc lập thẩm định đầy đủ. Ngoài ra, corpus
phải được gắn phiên bản và ngày hiệu lực để tránh sử dụng một quy định đã bị
sửa đổi hoặc hết hiệu lực. Cho đến khi hoàn thành hai yêu cầu này, kết quả
chỉ nên được xem là đánh giá kỹ thuật trên một benchmark nghiên cứu.

\subsection{Sử dụng có trách nhiệm}

Hệ thống là nguyên mẫu nghiên cứu, không phải công cụ tư vấn pháp lý và
không được dùng để tự động đưa ra quyết định ảnh hưởng đến quyền hoặc nghĩa
vụ của cá nhân. Một câu trả lời có citation vẫn có thể dựa trên path sai,
điều luật thiếu hoặc biểu diễn quy tắc không đầy đủ.

Nếu được phát triển thành ứng dụng hỗ trợ, hệ thống cần:

\begin{itemize}
    \item hiển thị rõ trạng thái \texttt{UNCERTAIN} và không ép thành kết luận;
    \item cho phép kiểm tra từng rule, article và hop của path;
    \item ghi lại phiên bản corpus, model, code và thời điểm truy vấn;
    \item cảnh báo rằng output cần được chuyên gia xác minh;
    \item bảo vệ dữ liệu cá nhân nếu sau này xử lý hồ sơ hoặc tình huống thật;
    \item duy trì cơ chế cập nhật và thu hồi tri thức pháp luật lỗi thời.
\end{itemize}

Các hạn chế trên xác định ranh giới cho mọi tuyên bố của nghiên cứu và là cơ
sở cho giai đoạn thực nghiệm tiếp theo.